% =====================================================================
%  Chapter 18 — Production Mechanisms
% =====================================================================

\chapter{Production Mechanisms}
\label{ch:production}

This chapter asks a single question for each mechanism: given that dark
matter exists now with the density it has, what had to have happened in the
early universe, and how much of that is explained rather than assumed?

\section{Thermal freeze-out}The mechanism of Chapter~\ref{ch:wimp}, stated as a
process rather than a candidate. A species in equilibrium with a bath
annihilates faster than the universe expands; as the expansion rate
outgrows the annihilation rate, the population freezes out and the
abundance is then diluted by the expansion.

What it explains: the density, in the sense that a standard-scale particle
with a standard-scale cross section produces it with no free parameter.

What it assumes: that the particle was in equilibrium, that the bath was the
Standard Model plasma, that nothing else interacted with it, and that the
temperature was standard. Each of these is a cosmological assumption, and
each is independently constrained.

What it gets wrong: the candidate.

\section{Thermal freeze-in}A species that never reaches equilibrium is produced
slowly by a coupling to the Standard Model, and the final abundance is
proportional to that coupling times the number of expansions. The result is
that the abundance is set by a parameter --- the coupling --- which is small
enough for the production to be slow.

What it explains: nothing about the density. The density is the coupling
re-expressed, with the physical scale removed.

What it assumes: the coupling, and therefore the initial conditions.

What it gets right: it makes the dark matter very weakly coupled, which is
consistent with everything the experiments have found and which is a
reasonable expectation for a species with a small interaction.

The assessment is uncomfortable but simple. Freeze-in is a mechanism that
converts the abundance question into a coupling question and then declines
to answer it. It is a legitimate part of the landscape and a poor place to
invest.

\section{Misalignment}A field with a vacuum expectation value that is displaced
from its minimum when the field-dominated era begins oscillates about the
minimum, and the coherent energy in the oscillation is converted into
particles when the field decays. The abundance depends on the initial
displacement angle $\theta_i$:

\[
\Omega_a h^2 \;\approx\; 0.15\,
\left(\frac{\theta_i}{0.01}\right)^{2}
\left(\frac{f_a}{10^{12}\ \GeV}\right)^{1.165}.
\]

The scale $f_a$ is the symmetry-breaking scale, which is a parameter of the
model, and the angle $\theta_i$ is a parameter of the history.

What it explains: the density, for a natural angle of order one and a
natural symmetry-breaking scale. The density is not predicted, but it is
reachable with inputs that were chosen for other reasons --- $f_a$ is
constrained by the Peccei--Quinn solution to the strong CP problem, and the
angle of order one is not a special value.

What it assumes: that the initial angle was not fine-tuned. For the QCD
axion this is defensible; for a generic ultralight field with a much larger
scale it is not, as Chapter~\ref{ch:ultralight} noted.

\section{Gravitational production}A field that is not in thermal equilibrium because
its coupling is negligible can still be produced by the expansion of the
universe itself. The classic result is that a very light field can be
produced abundantly as the inflaton's fluctuations are stretched beyond the
horizon, with a production efficiency that peaks at a mass set by the
inflationary energy scale.

What it explains: the abundance, with no free parameter other than the
inflationary history.

What it assumes: the inflationary energy scale and the inflaton properties,
which are constrained independently by the primordial power spectrum, the
tensor-to-scalar ratio, and the non-observation of spectral features.

What it gets right: it is the natural production mechanism for a very light
field, which is the class whose production is otherwise most awkward.

\section{Hawking production and superheavy relics}Black holes formed in the early universe
emit Hawking radiation, and any particle species light enough relative to
the black hole temperature is produced in an abundant flux. The abundance is
determined by the number density of black holes, which is determined by the
primordial power spectrum, and the mass distribution of the relics follows
the mass function of the black holes.

What it explains: the abundance, for a specified inflationary spectrum.

What it assumes: the primordial power spectrum at small scales, which is
constrained by the CMB spectral distortions and by the absence of features
in the CMB spectrum.

\section{A summary comparison}Five mechanisms, and the comparison that matters is not
which one is most popular but which ones can be falsified by the density
constraint itself.

\begin{table}[H]
\centering
\caption{Production mechanisms compared. The third column is the decisive
one: a mechanism whose abundance depends on a predicted quantity can be
wrong in a way that the abundance constraint detects.}
\label{tab:production}
\begin{tabular}{@{}>{\RaggedRight\arraybackslash}p{0.24\textwidth}p{0.30\textwidth}p{0.32\textwidth}@{}}
\toprule
\textbf{Mechanism} & \textbf{Abundance set by} & \textbf{Falsifiable by the density?} \\
\midrule
Freeze-out & $\langle\sigma v\rangle$ & Yes, and it has been \\
Freeze-in & A coupling parameter & No \\
Misalignment & $\theta_i$ and $f_a$ & Only indirectly \\
Gravitational & Inflationary history & Yes, if the spectrum is measured \\
Hawking & Primordial power spectrum & Yes, if the spectrum is measured \\
\bottomrule
\end{tabular}
\end{table}

The pattern in \tabref{tab:production} is a guide to the strategy of
Chapter~\ref{ch:portfolio}. The mechanisms that are falsifiable by the
density constraint are also the mechanisms for which a measurement of the
primordial spectrum, or of the CMB spectral distortions, is a dark matter
experiment. That is a much cheaper test than any detector, and it is
under-used.

\begin{keytakeaways}
\item Each production mechanism has a different profile: some predict the
abundance from quantities that can be measured independently, some only
parameterise it.
\item Gravitational and Hawking production turn a measurement of the
primordial power spectrum into a dark matter measurement --- one of the
cheapest tests available and one of the least used.
\item Misalignment is defensible for the QCD axion, where the scale is fixed
by another problem, and is not defensible as a generic explanation.
\item Freeze-in converts the density question into a coupling question and
leaves it there.
\end{keytakeaways}
